% Transcription — fonds Alexandre Grothendieck, shelfmark 12, batch 5 (pages 81-100).
% Established from the facsimile archives/batches/12/batch-05.pdf (a mirror of
% https://grothendieck.umontpellier.fr/12.pdf), per the skill
% transcribe-grothendieck. The batch file carries no cover sheet: PDF page N
% is page 80+N, as the Montpellier footer printed on each PDF page confirms.
% The archivists' pencil numbers bottom-left agree with that position on every
% leaf (81 on the first, 100 on the last), so \page{} takes them as they are.
%
% Pass: Opus 5.5 (claude-opus-5-5), 2026-09-23 — first pass, unchecked against the pages by a human.
%
% What the batch is. Loose leaves in his hand, a numbered run, and a
% typescript of his own, in that order.
%   81      A note on its own: C a proper smooth curve over F_q, A an abelian
%           scheme over C; H^i(C,A) = 0 for i >= 2 by Weil, H^1(C,A) = 0
%           "sans doute" by Tate; the localisation sequence gives
%           H^1(U,A) = H^2_x(C,A) = H^1(K_x,A), dual of H^0(K_x,A*), a finite
%           group. Closes on whether this generalises to X regular minus a
%           regular divisor.
%   82, 84  Skipped. Typed leaves of an EGA IV draft (a proposal to restate
%           Théorème 7.5.1; Remarques, Définition 5 and Lemme 6 on
%           schematically dominant families), cancelled by long diagonal
%           strokes and used as scrap: they are the backs of 81 and 83. Not
%           the mathematics of this folder, and nothing on them is in his
%           hand beyond inked Greek letters.
%   83      « Calculs sur T_l(G_m) = Q_l(1) et G_m »: X of finite type over Z,
%           H^i(X, Q_l(1)) for i = 0..3 and i >= 4 in terms of units, Pic and
%           the (co)Brauer group; Artin's conjecture that Br is finite for X
%           proper over Z; how H^3 behaves on passing to an open.
%   85-88   His § 17, « Formes positives »: for a semi-simple pure motive M
%           of weight p over an algebraically closed k, the Q-space L^s(M)
%           of symmetric (p even) / alternating (p odd) forms
%           M x M -> Q(-p); a cone L^+(M) of positive forms, its interior the
%           ample forms; eight axioms (i)-(viii); consequences: compatibility
%           with Tate twists, positive definiteness on Gamma_k M in weight 0,
%           reduction to isotypic pieces I (x)_A M.
%   89-94   His § 18, « Dictionnaire Fonctions L – Cohomologie à action
%           galoisienne »: the ring R(N;k) of k[N] = k[t]-modules finite
%           projective over k; L(V,f) = 1/det(1 - ft) as a lambda-ring map to
%           1 + k[[t]]^+, injective over a field, and its logarithmic
%           derivative sum (Tr f^n) t^(n-1); Lefschetz–Weil–Grothendieck:
%           L(X,F) determines chi(F) in R(N,Q_l), and under Weil its weight
%           pieces h^i(F); birational invariants from L modulo lower weights
%           (p. 92, with a struck-out passage on a fibration f : X -> Y), the
%           two top compactly supported groups (p. 93), and a programme
%           (p. 94) relative to a morphism f : X -> Y. The margin of p. 92
%           refers to his « N° 12 ».
%   95      Cover in his hand, « Poids et niveaux d'espaces de cohomologie
%           l-adiques (base var. quelconque) ». No \page; its words are the
%           section heading.
%   96-97   His typescript « Filtration et poids des espaces de cohomologie
%           des variétés algébriques par A. GROTHENDIECK », § 0 Introduction,
%           with corrections in his hand; the l's are inked by hand. It says
%           the idea of the weight arguments « remonte à 1964 » and cites
%           SGA 7 and Deligne's Hodge theory; nothing on the leaf is dated.
%   98      His handwritten plan of the same paper: §§ 0-9 and a list of
%           references with the brackets left empty.
%   99      A typed leaf in English, « 2. Formalism of Algebraic Cohomology
%           Classes », the cycle map rho_X : C*(X) -> H*_l(X), formulas and
%           a few interlinear corrections inked by hand. The leaf does not
%           name its author and the hand of the corrections is not certainly
%           his; it is transcribed because it is the mathematics the plan of
%           p. 98 refers to (« Kleiman / conj. standard »), and a \note says
%           so.
%   100     Calculations in his hand on a counterexample: T two crossing
%           lines in P^2, T' two points a, b, X = P^2 - T', Y = T - T',
%           U = P^2 - T, and the long exact sequences of compactly supported
%           cohomology; « inj ? ».
%
% Notation fixed for the batch. Blackboard letters for Q, Z, F_q, P, G_m, N;
% Q_\ell for his Q_l. The script letter of p. 85 onwards (L with a bar) is
% \mathcal{L}, with superscripts s and +; the script B of p. 86 is
% \mathcal{B}. His compactly supported cohomology H_! is H_{!}. The letter he
% overwrites heavily inside k[·] and R(·;k) on pp. 89-91 is read \mathbb{N}
% (the monoid algebra k[N] = k[t]); a \note says so once. His Z(X,F) of
% p. 93, drawn as a figure-3, is set \zeta. His underlining in prose is
% \emph.
%
% Typescript x-outs. On pp. 96-97 the typist's own overstrikes (xxx over a
% false start) are transcribed as \struck{} only where they carry a phrase
% (e.g. « qui remonte à 1964 », « Pour terminer cette introduction » typed in
% the wrong place); single-letter overtypes are silently resolved.
%
% Legibility. Pages 96-99 read through (typescript); 83, 85, 88, 98 and 100
% are written fair; 86-87, 90-94 are a fast draft whose formulas carry and
% whose connective prose often does not; the diagonal margins of 92 and 93
% are largely lost.
\documentclass[11pt,a4paper]{article}
% Le préambule commun à toutes les transcriptions du fonds Grothendieck.
%
% Il tient en une idée : **l'incertitude se compose, elle ne se résout pas.**
% Une transcription qui lisse un mot douteux a détruit la seule chose pour
% laquelle on la faisait — un lecteur doit pouvoir savoir, à chaque endroit,
% ce qui était sur le feuillet et ce qui a été ajouté. Les six macros qui
% suivent sont donc l'appareil critique, et non de la décoration.
%
% Les mêmes commandes sont reconnues par `scripts/render.mjs`, qui produit la
% vue de lecture du site. Une macro ajoutée ici doit l'être aussi là-bas,
% faute de quoi le rendu échouera bruyamment — ce qui est voulu : un rendu
% silencieusement faux est pire qu'un rendu absent.

\ProvidesPackage{grothendieck}[2026/08/08 Transcriptions du fonds Grothendieck]

\usepackage{amsmath,amssymb,amsthm}
\usepackage{mathrsfs}
\usepackage{stmaryrd}
% Les diagrammes commutatifs sont partout dans ce fonds : c'est la langue
% même de la géométrie algébrique de Grothendieck, et une transcription qui
% les rendrait en prose ne transcrirait rien.
\usepackage{tikz-cd}
% Français par défaut — c'est la langue des feuillets — mais l'anglais est
% chargé aussi : la traduction et le résumé sont des documents anglais, et la
% typographie française (espaces avant les deux-points) y serait une faute.
% Ces documents basculent par \selectlanguage{english} en tête de corps.
\usepackage[english,french]{babel}
\usepackage{fontspec}
\usepackage{geometry}
\geometry{a4paper,margin=25mm,marginparwidth=28mm}
\usepackage{marginnote}
\usepackage{xcolor}
% ulem, not soul. `soul`'s \st and \ul are built on a character-by-character
% scan that XeLaTeX's font machinery breaks: the document fails with an
% undefined control sequence rather than degrading. ulem does the same two jobs
% and is Unicode-engine safe. `normalem` stops it hijacking \emph, which the
% transcriptions use for Grothendieck's own emphasis.
\usepackage[normalem]{ulem}
\usepackage{hyperref}
\hypersetup{colorlinks=true,linkcolor=black,urlcolor=black,pdfborder={0 0 0}}

% --- Métadonnées du lot -----------------------------------------------------
% Reprises telles quelles par le script de rendu, qui les affiche en tête de
% la vue de lecture. Elles fixent ce que le fichier prétend couvrir : sans
% elles, une transcription flotte sans point d'attache dans un fonds de seize
% mille feuillets.

\newcommand{\folder}[1]{\gdef\grothFolder{#1}}
\newcommand{\batch}[1]{\gdef\grothBatch{#1}}
\newcommand{\leaves}[2]{\gdef\grothFirst{#1}\gdef\grothLast{#2}}
\newcommand{\foldertitle}[1]{\gdef\grothFoldertitle{#1}}
\gdef\grothFolder{?}\gdef\grothBatch{?}\gdef\grothFirst{?}\gdef\grothLast{?}\gdef\grothFoldertitle{}

% --- Le résumé --------------------------------------------------------------
%
% Il ouvre la lecture modernisée, sous le titre « Résumé », et il est composé
% en petit. Ce n'est pas un ornement : c'est le seul passage du document qui
% ne soit pas des mathématiques, et un lecteur doit pouvoir le reconnaître
% avant de l'avoir lu — puis le sauter s'il connaît déjà le sujet, ou s'y
% arrêter s'il ne le connaît pas. Le filet qui le suit marque la reprise du
% corps.

\newenvironment{resume}
  {\small\setlength{\parskip}{0.45em}}
  {\par\medskip\hrule\medskip}

% --- Mots-clés ---------------------------------------------------------------
%
% En anglais, en clôture du résumé de la lecture modernisée : le vocabulaire
% moderne sous lequel chercher ce que les feuillets construisent. Le manifeste
% du site (`npm run manifest`) les extrait pour étiqueter la cote entière —
% cette ligne est la source unique des tags, il n'y a pas de fichier à côté.
\newcommand{\keywords}[1]{\par\smallskip\noindent
  {\footnotesize\textsc{Keywords} --- \textit{#1}}\par}

% --- L'appareil critique ----------------------------------------------------

% Le numéro de feuillet, en marge. C'est l'ancre qui permet de régler un
% désaccord sur une formule en regardant *un* feuillet plutôt qu'une chemise
% de six cents. Le site s'en sert aussi pour tourner les pages du fac-similé
% quand on fait défiler la transcription.
\newcommand{\leaf}[1]{%
  \marginnote{\footnotesize\color{gray!70}\textsf{#1}}%
  \label{leaf:#1}\ignorespaces}

% Illisible : on ne devine pas. Un mot inventé qui se lit comme les autres est
% le pire résultat possible.
\newcommand{\ill}{\textcolor{red!60!black}{[\,\ldots\,]}}

% Lecture douteuse : le mot est donné, et signalé comme incertain.
\newcommand{\uncertain}[1]{\uline{#1}}

% Ajout de l'éditeur — une lettre manquante, un mot sous-entendu.
\newcommand{\add}[1]{\textcolor{blue!50!black}{[#1]}}

% Biffé par Grothendieck lui-même. Ce qu'il a rayé fait partie du document :
% une rature dit souvent où la pensée a changé de direction.
\newcommand{\struck}[1]{\sout{#1}}

% Note du transcripteur, sur l'état matériel du feuillet.
\newcommand{\note}[1]{\par\smallskip{\small\color{gray!60!black}\textit{[#1]}}\par\smallskip}

% Note marginale de Grothendieck, distincte du corps du texte.
\newcommand{\marginal}[1]{\par\smallskip\noindent\hspace{1em}%
  {\small\color{gray!60!black}#1}\par\smallskip}

% --- Titre ------------------------------------------------------------------

\AtBeginDocument{%
  \begin{center}
    {\large\bfseries Fonds Alexandre Grothendieck --- cote n\textsuperscript{o}~\grothFolder}\\[2pt]
    {\itshape\grothFoldertitle}\\[4pt]
    {\small feuillets \grothFirst--\grothLast\ (lot \grothBatch)}\\[2pt]
    {\footnotesize\color{gray!60!black}Transcription établie d'après le fac-similé numérisé
     par l'Université de Montpellier.\\
     Les lectures incertaines sont soulignées, les illisibles notées [\,\ldots\,],
     les ajouts entre crochets.}
  \end{center}
  \vspace{1em}}

\endinput


\folder{12}
\batch{5}
\pages{81}{100}
\dating{1964-[vers 1971]}
\watermark{Édition de démonstration}
\foldertitle{Généralités (sous-groupes de Galois motiviques) (1964 ?) : notes manuscrites (s.d.), tapuscrit (s.d.)}

\begin{document}

\page{81}
$C$ courbe lisse propre sur $\mathbb{F}_{q}$
\note{« propre » est ajouté en interligne}

$A$ schéma abélien sur $C$.

Des conjectures de Weil, résulte que
\[
  H^{i}(C, A) = 0 \quad \text{si } i \geqslant 2
\]
et sans doute, utilisant Tate, on doit avoir
\[
  H^{1}(C, A) = 0
\]
Soit $x$ un pt fermé de $C$, $U = C - \{x\}$, alors

\begin{tikzcd}[column sep=small]
  \cdots \arrow[r] & H^{1}(C, A) \arrow[r] \arrow[d, no head, "\Vert" description] & H^{1}(U, A) \arrow[r] & H^{2}_{x}(C, A) \arrow[r] & H^{2}(C, A) \arrow[d, no head, "\Vert" description] \\
  & 0 & & & 0
\end{tikzcd}

\note{le premier terme de la suite est biffé et illisible ; il est remplacé ici par $\cdots$}
d'où
\[
  H^{1}(U, A) \simeq H^{2}_{x}(C, A)
\]
Or par Tate,
\[
  H^{2}_{x}(C, A) \simeq H^{1}(K_{x}, A)
\]
est \emph{dual} de $H^{0}(K_{x}, A^{*})$
\note{les deux signes $\simeq$ sont surchargés}
\struck{\ill{} un \uncertain{groupe fini} \ill{} \uncertain{c'est donc} \ill{}
\uncertain{extension} \ill{} \uncertain{$p$-groupe}, \ill{}}
\note{deux lignes biffées d'un double trait}
groupe fini ! OK.

Peut-on généraliser au cas d'une variété $X$ régulière dont on enlève un
diviseur régulier ---

\page{83}
\emph{Calculs} sur $T_{\ell}(\mathbb{G}_{m})$ $= \mathbb{Q}_{\ell}(1)$
et $\mathbb{G}_{m}$.
\note{« $= \mathbb{Q}_{\ell}(1)$ » est ajouté au-dessus de la ligne}

$X$ de type fini sur $\mathbb{Z}$
\[
\begin{array}{lll}
  H^{0}(X, \mathbb{Q}_{\ell}(1)) & = 0 & \\[1ex]
  H^{1}(X, \mathbb{Q}_{\ell}(1)) & = H^{0}(X, \mathcal{O}_{X}^{*}) \otimes_{\mathbb{Z}} \mathbb{Q}_{\ell} & \\[1ex]
  H^{2}(X, \mathbb{Q}_{\ell}(1)) & \text{extension de } T_{\ell}(H^{2}(X, \mathbb{G}_{m})) & \text{par } H^{1}(X, \mathcal{O}_{X}^{*}) \otimes \mathbb{Q}_{\ell}
\end{array}
\]
\note{dans la première ligne, un argument biffé après $\mathbb{Q}_{\ell}(1)$ ;
les termes $H^{0}(X, \mathcal{O}_{X}^{*}) \otimes \mathbb{Q}_{\ell}$,
$T_{\ell}(H^{2}(X, \mathbb{G}_{m}))$ et $H^{1}(X, \mathcal{O}_{X}^{*}) \otimes
\mathbb{Q}_{\ell}$ sont entourés}
\marginal{$\sim$ \uncertain{coBrauer}}
\[
  H^{3}(X, \mathbb{Q}_{\ell}(1)) \simeq T_{\ell}(H^{3}(X, \mathbb{G}_{m}))
  \ \text{est ce qu'il peut}
\]
[p.\ ex.\ $\mathbb{Z}_{\ell}$ si $X$ \uncertain{modèle} \ill{} d'un corps de
fonctions d'une variable sur $\mathbb{F}_{p}$]
\[
  H^{i}(X, \mathbb{Q}_{\ell}(1)) = 0 \quad \text{si } i \geqslant 4 \ ??
\]
\[
  \Vert
\]
\[
  T_{\ell}(H^{i}(X, \mathbb{G}_{m}))
\]
On conjecture (Artin) que $H^{2}(X, \mathbb{G}_{m})$ = \uncertain{fini} si $X$
propre sur $\mathbb{Z}$. \struck{Mais si} Vrai si $X$ courbe projective lisse
sur $\mathbb{F}_{p}$, mais \uncertain{faux quand on} lui enlève des pts -- le
rang du $H^{2}$ augmente d'autant
\note{« Vrai » est écrit au-dessus de « Mais si », biffé}

$H^{3}(X, \mathbb{Q}_{\ell}(1))$ : quelle est sa taille [il décroît en
passant à un ouvert de $X$]

\section{17 Formes positives}
\note{titre souligné, de sa main, précédé du numéro 17 qu'il donne au
paragraphe}

\page{85}
Le théorème étant géométrique, nous nous plaçons sur un corps alg.\ clos
$k$. Nous ne considérons que les motifs ½ \emph{simples} sur $k$,
purs, de poids $p$. Soit $M$ un tel motif, considérons les
formes
\[
  M \times M \longrightarrow \mathbb{Q}(-p)
\]
i.e.\ les éléments de
\[
  \Gamma_{k}(\operatorname{Hom}(M \otimes M, \mathbb{Q}(-p))) = \mathcal{L}(M)
\]
où le \struck{\ill{}} $\operatorname{Hom}$ est bien de poids $0$.
\note{« purs » est ajouté en interligne}

Elles forment un vectoriel de \struck{poids} dim.\ finie sur
$\mathbb{Q}$. On considère \uncertain{le sous-espace} \ill{} dans
\note{« dim.\ finie » est écrit au-dessus de « poids », biffé}
formé des formes \emph{symétriques} si $p$ pair, \emph{alternées} si $p$
impair. Soit \uncertain{ce que nous notons} $\mathcal{L}^{s}(M)$. Noter les
hom.\ canoniques
\[
  \mathcal{L}^{s}(M) \otimes \mathcal{L}^{s}(N) \longrightarrow \mathcal{L}^{s}(M \otimes N)
\]

\page{86}
par produit tensoriel des formes [N.B.\ le produit tensoriel de deux formes
\struck{symétriques, ou de} de même type de symétrie est symétrique -- de deux
types opposés de symétrie est antisymétrique]. Bien entendu, ils ne sont pas
en général bijectifs (mais du moins injectifs).

On a aussi
\[
  \mathcal{L}^{s}(M' \times M'') \simeq \mathcal{L}^{s}(M') \times \mathcal{L}^{s}(M'') \times \mathcal{B}(M', M'')
\]
où
\[
  \mathcal{B}(M', M'') = \operatorname{Hom}(M' \otimes M'', \mathbb{Q}(p)) .
\]
\note{$\mathbb{Q}(p)$ sic, là où la p.~85 écrit $\mathbb{Q}(-p)$}

Ceci posé, on \struck{\ill{}} \ill{} $\mathcal{L}^{s}(M)$ un cône
\uncertain{convexe} \struck{\uncertain{fermé}} $\mathcal{L}^{+}(M)$.
\note{« convexe » est ajouté au-dessus d'un mot biffé ; au-dessus de la
ligne, un mot biffé et « des », puis une lettre isolée, peut-être $H$}

On a les propriétés suivantes :

\marginal{les formes $\in \mathcal{L}^{+}(M)$ sont dites \uncertain{positives},
celles dans l'intérieur sont dites amples}
\begin{enumerate}
  \item[(i)] $\mathcal{L}^{+}(M)$ a un intérieur non vide.
  \item[(ii)] Un élément de $\mathcal{L}^{+}(M)$ est dans l'int.\ de
  $\mathcal{L}^{+}(M)$ ssi il est non dég.
  \item[(iii)] Pour $M \to N$, $\mathcal{L}^{+}(N)$ s'envoie dans
  $\mathcal{L}^{+}(M)$.
  \item[(iv)] \struck{\ill{}} Si $M \to N$ injectif, \struck{$\mathcal{L}^{+}(N)$}
  l'intérieur de $\mathcal{L}^{+}(N)$ s'envoie dans l'intérieur de
  $\mathcal{L}^{+}(M)$ (formes amples $\longrightarrow$ formes amples)
\end{enumerate}

\page{87}
\struck{\ill{}} i.e.\ la forme induite sur $M$ par une forme ample est
non dégénérée, donc
\[
  N = M \oplus M^{\perp} ,
\]
et la forme sur $N$ est la somme des formes induites sur $M$, $M^{\perp}$.
\note{« par une forme ample » est ajouté au-dessus de la ligne}

\begin{enumerate}
  \item[(v)] La somme directe de deux formes \struck{positives ou}
  pseudo-amples (amples) est \uncertain{item}.
  \item[(vi)] Le produit tensoriel de deux formes amples est ample.
  \item[(vii)] Si $p = 0$, de sorte que \uncertain{essentiellement} $M$ est
  défini par un vectoriel de dim.\ finie sur $\mathbb{Q}$, alors les formes
  amples sont les formes \emph{définies positives}. \struck{\ill{} $M(p)$,
  $M$ un \uncertain{vectoriel} \ill{} de dim.\ finie sur $\mathbb{Q}$ ;
  utilisant l'isom.\ canon.}
  \item[(viii)] Cas de $\mathbb{Q}(p)$ : on a $\mathcal{L}^{s}(\mathbb{Q}(p))
  = \mathbb{Q}$, et on veut que les formes amples correspondent aux nbs
  \emph{strictement} positifs.
\end{enumerate}
\note{à la fin de (vii), un passage de trois lignes est biffé d'un trait
ondulé et encadré ; « strictement » est ajouté au-dessus de la ligne en
(viii)}

\page{88}
De (vi) et (viii) résulte que pour tout motif $M$ et entier $n$,
l'isomorphisme
\[
  \mathcal{L}^{s}(M) \simeq \mathcal{L}^{s}(M(n))
\]
est compatible avec les cônes amples.

De (iv) et \uncertain{(vi)} résulte que si $p = 0$, une forme ample sur $M$
induit une forme \emph{définie positive}
\note{le renvoi attendu serait plutôt (vii) ; on transcrit ce qui paraît
écrit}
\[
  \Gamma_{k} M \times \Gamma_{k} M \longrightarrow \Gamma_{k} \mathbb{Q}(0) = \mathbb{Q}
\]

\uncertain{En décomposant} $M$ (de poids $p$) en composantes isotypiques
$M_{\alpha}$, il vient
\[
  \mathcal{L}^{s}(M) \simeq \bigoplus_{\alpha} \mathcal{L}^{s}(M_{\alpha})
\]
\note{le signe de somme directe est surchargé}
donc pour étudier les formes positives, il suffit en principe de les étudier
sur les morceaux isotypiques, qui sont de la forme
\[
  I \otimes_{A} M
\]
où $I$ est un motif simple, $A$ son anneau des endomorphismes, qui est un
corps, et $M$ un \emph{vectoriel} sur $A$

\section{18 Dictionnaire Fonctions L -- Cohomologie à action galoisienne}
\note{« Dictionnaire » doublement souligné ; le numéro 18 est encerclé, en
tête du feuillet ; « à action galoisienne » est écrit sous la fin du titre}

\page{89}
Soit $k$ un \struck{corps} anneau

\struck{$A$ la catégorie} $\mathcal{C}$ la catégorie des \struck{modules} de
\uncertain{longueur} \struck{$\operatorname{Mod}(\mathbb{Z})$ la catégorie}
finie sur $k[t]$ = modules $V$ proj.\ de \struck{vectoriels}
\struck{\uncertain{de dim.}} finis sur $k$ munis d'un endomorphisme $f$ =
$k[\mathbb{N}]$-modules, projectifs \struck{de type} finis sur $k$.
\note{« proj.\ de » et « projectifs » sont ajoutés au-dessus de la ligne ; la
lettre entre crochets, dans $k[\,\cdot\,]$ ici et dans $R(\,\cdot\,; k)$ plus
bas, est fortement surchargée et se lit $\mathbb{N}$ : $k[\mathbb{N}] = k[t]$}

Cette catégorie donne lieu à un \uncertain{$\lambda$-anneau}, noté
$R(\mathbb{N}; k)$, qui est aussi, si $k$ est un corps, le
\struck{\ill{}} \uncertain{$\mathbb{Z}$-libre} engendré par les classes des
$k[\mathbb{N}]$-modules simples finis sur $k$, ou encore les
$k[t]$-modules simples, ou encore les polynômes unitaires
irréductibles $P \in k[t]$, en envoyant : $P$ en \ill{} \ill{}. À tout tel
$P$, on associe : $k[t]/P\,k[t]$.
\note{« si $k$ est un corps », « finis sur $k$ », « unitaires » sont ajoutés
en interligne}
\marginal{\ill{}}
$(V, f)$, dont le polynôme car.\ est précisément $P$,
\[
  P(t) = \det(t \operatorname{id} - f)
\]
[vrai pour tt $P$, irréd.\ ou non !]. \uncertain{Il nous} sera plus commode
de remplacer $P(t)$ par
\[
  Q(t) = t^{d} P(\tfrac{1}{t}) = \det(1 - tf)
\]
qui est un polynôme quelconque à coeff.\ cst égal à $1$.

\page{90}
Ainsi, on obtient via la \uncertain{fonction} car.\ modifiée un hom
\[
  (V, f) \rightsquigarrow \frac{1}{\det(\operatorname{id}_{V} - f t)} = L(V, f)
\]
\[
  R(\mathbb{N}, k) \longrightarrow 1 + k[[t]]^{+}
\]
[compatible avec \struck{\ill{}} \uncertain{l'anti}-« additif »
$\varphi(t) \mapsto \dfrac{1}{\varphi(t)}$ dans $1 + k[[t]]^{+}$] qui est un
hom.\ de $\lambda$-anneaux,
\marginal{en fait, \ill{} les \emph{fonctions rationnelles} de valeur $1$ en
$0$}
\note{la note marginale, encadrée, est reliée par une flèche à
$R(\mathbb{N}, k)$}

Lorsque $k$ est \struck{un} corps, l'hom.\ précédent est injectif. En effet,
si \struck{\ill{}}, la connaissance de la fonction rationnelle à
numérateur $1$ dans $k[t]$ \ill{} \uncertain{facilement} \ill{}
[$k[t]$ étant factoriel] l'élément
\[
  \sum_{\substack{P \text{ pl.} \\ \text{unitaire}}} c_{P} (V_{P}, f_{P})
\]
\struck{dont} par
\[
  L = \prod_{Q} Q^{c_{P}}
\]
[$Q$ et $P$ associés de la façon précisée plus haut].

Notons qu'on a un homomorphisme naturel \struck{\ill{}}
\[
  1 + k[[t]]^{+} \longrightarrow k[[t]]
\]
\[
  \varphi(t) \rightsquigarrow \frac{\varphi'(t)}{\varphi(t)}
\]
qui correspond à :
\[
  \Lambda(V, f)(t) = \frac{L(V, f)'(t)}{L(V, f)(t)} = [\log L(V, f)]' = \frac{1}{t} \sum_{1}^{\infty} (\operatorname{Tr} f^{n})\, t^{n}
\]
\struck{En car.\ $0$, on} Si $k$ de car.\ $0$, i.e.\ $k \supset
\mathbb{Q}$, on trouve inversement $L$ par
$\exp\bigl(\int_{0}^{t} \Lambda(V, f)(t)\bigr)$. D\ill{}
\marginal{un corps de car.\ $0$}
\[
  \Lambda : R(\mathbb{N}, k) \longrightarrow k[[t]]
\]
est alors injectif si $k$ corps \struck{($k$ corps}.
\note{« un corps de car.\ $0$ », dans la marge gauche, est relié par une
flèche à la dernière ligne}

\page{91}
Si $k$ n'est pas de car.\ $0$, on ne peut en général aller en sens inverse,
et $R(\mathbb{N}, k) \xrightarrow{\ \Lambda\ } k[[t]]$ n'est \struck{en
général} pas injectif. [$k$ corps de car.\ $> 0$]

On voit \struck{que} par la formule de Lefschetz--Weil--Grothendieck que la
connaissance de la fonction $L$ de $(X, F)$ \struck{\ill{}} à coeff.
dans $F$, sur $\mathbb{F}_{p}$, \emph{équivaut} à celle de l'élément
\[
  \chi(F) = \sum (-1)^{i} \operatorname{cl} H^{i}_{!}(\overline{X}, \overline{F}) \quad \text{de } R(\mathbb{N}, \mathbb{Q}_{\ell})
\]
[on considère cette fonction $L$ comme l'image de $\chi(F)$ dans $1 +
\mathbb{Q}_{\ell}[[t]]^{+}$]. Lorsqu'on admet les conjectures de Weil, et $F$
\emph{\uncertain{vient d'un motif}}, $\chi(F)$ se décompose suivant les poids
en
\[
  \sum (-1)^{i} \underline{h}^{i}(F) , \qquad \underline{h}^{i}(F) \in R(\mathbb{N}, \mathbb{Q}_{\ell}) ,
\]
et la \struck{connaissance} connaissance de $L(X, F)(t)$ \emph{équivaut} à
celle des $\underline{h}^{i}(F)$ : les $\underline{h}^{i}(F)$ correspondent à
la décomposition de $L(X, F)$

\page{92}
\struck{facteurs} en groupant suivant les valeurs absolues des zéros et
pôles… Lorsque $X$ est complet et lisse, et $F$ localement constant, de
poids $p$, on a
\[
  \underline{h}^{i}(F) = H^{i}(\overline{X}, \overline{F})
\]
comme $\mathbb{N}$-modules, donc on a la connaissance de $L(X, F)$ équivaut à
celle des $H^{i}(\overline{X}, \overline{F})$ et l'opération de Frobenius
dans ceux-ci.
\note{« et $F$ localement constant, de poids $p$ » est ajouté au-dessus de la
ligne}

\struck{Soit $f : X \to Y$ un morphisme [\ill{} propre] avec $X$, $Y$
\ill{} de type fini sur $\mathbb{F}_{p}$. Considérons les ouverts denses $U$
de $Y$, et les $L(f^{-1}(U), F)$. Pour $U$ variable, ces [pour $X$ de dim.
$n$] fonctions rationnelles}
\note{tout ce passage est encadré et annulé de hachures obliques ; il porte à
gauche le diagramme suivant, où $K$ = corps des fonctions de $Y$}

\begin{tikzcd}
  X \arrow[r, no head] \arrow[d] & X_{K} \arrow[d, no head] \\
  Y \arrow[r, no head] & K
\end{tikzcd}

Considérons, d'autre part, les ouverts \struck{denses} $U$ de $X$, tels
que $X - U$ de dim.\ $< n$, et les $L(U, F|U)$, qui ne diffèrent
\struck{que} (pour $F$ \uncertain{pur} de poids $p$) que par des fonctions
rationnelles de poids $\leqslant 2n - 2 + p$. Donc, à priori la connaissance
de l'invariant birationnel $L(X, F)$ mod.\ fonctions rat.\ de poids
$\leqslant 2n - 2 + p$,
\marginal{\uncertain{le poids} \ill{} \uncertain{filtration} \ill{} cf.
\uncertain{N° 12}}

\page{93}
\emph{équivaut} à celle de $\operatorname{cl} H^{2n}_{!}(X, F) -
\operatorname{cl} H^{2n-1}_{!}(X, F)$ \ill{} la \ill{} chose. On peut
supposer $X$ [\uncertain{affine} et] lisse sur $\mathbb{F}_{p}$ \uncertain{connexe}
de dim.\ $n$, alors $H^{2n}_{!}(X, F)$ est dual de
\[
  H^{0}(X, \operatorname{Hom}(F, \mathbb{Q}_{\ell}(n))) = H^{0}(X, \check{F})(n) ,
\]
\uncertain{donc} \ill{} des \uncertain{données} \ill{}. D'autre part,
$H^{2n-1}_{!}(X, F)$ est dual de
\[
  H^{1}(X, \operatorname{Hom}(F, \mathbb{Q}_{\ell}(n))) ,
\]
et \uncertain{admet} en général des poids $2n-1$ \emph{et} des poids $2n-2$.
\note{au-dessus de « de dim.\ $n$ », un mot biffé et un autre ajouté, lu
« connexe » sous réserve}
\marginal{\ill{} \uncertain{filtration} \ill{} \uncertain{poids} \ill{}
\uncertain{niveau} \ill{}}

\struck{Une analyse} Si $F = \mathbb{Q}_{\ell}(0)$, \struck{\ill{}} \ill{}
\struck{\ill{}} \ill{} des fonctions $\zeta$ habituelles, la connaissance
« birationnelle » grossière de $\zeta$ est \ill{} équivalente : celle de
$\pi_{0}(\overline{X})$ comme \ill{} \uncertain{fondamentale}, et celle de
Albanese de $X$ sur $\mathbb{F}_{p}$ ---
\note{la lettre qu'il trace en forme de 3 pour la fonction zêta est transcrite
$\zeta$}
\marginal{modulo conjectures de Tate}

Une analyse plus serrée \struck{des \ill{}} (plus \uncertain{fine} que par le
seul poids, mais en utilisant aussi le niveau) permet de \uncertain{récupérer}
plus \ill{} la connaissance de $\zeta(X, F)$ (disons, en prenant $F =
\mathbb{Q}_{\ell}$) modulo les $\zeta(Y)$ \struck{\ill{}} (dim.\ $Y < n$)
\struck{[poids $\leqslant p$]}

\page{94}
savoir les « invariants birationnels fondamentaux » relativement à
$\mathbb{F}_{p}$, envisagés dans les notes sur les \emph{Motifs}.

Il semblerait opportun d'étudier \struck{\ill{}} le cas des
coefficients $F$ généraux et de voir si, pour un morphisme $f : X \to Y$,
disons lisse propre \struck{\ill{}}, $F$ lct.\ constant sur $X$, la
connaissance \struck{des} de \struck{$H$}
\note{le paragraphe est marqué d'un double trait vertical dans la marge
gauche}
\[
  L(X, F) \ \text{modulo} \ L(f^{-1}(Y'), F) ,
\]
pour $Y'$ variant dans $Y$, permet de récupérer les invariants birationnels
fondamentaux de $(X_{K}, F|X_{K})/K$ [$K$ corps des fonctions de $Y$].
\emph{Dans quelle mesure faut-il faire les hypothèses} $f$ \emph{lisse
propre}, et $F$ \emph{lct.\ constant} pour obtenir des énoncés de cette
nature ??

\section{Poids et niveaux d'espaces de cohomologie $\ell$-adiques (base var.\ quelconque)}
\note{ces mots sont de sa main, seuls sur le feuillet 95, qui sert de
chemise aux pages suivantes}

\subsection{Filtration et poids des espaces de cohomologie des variétés algébriques}
\note{titre dactylographié, suivi de « par A.\ GROTHENDIECK » ; les
feuillets 96 et 97 sont un tapuscrit, corrigé de sa main, où les $\ell$ sont
ajoutés à l'encre}

\page{96}
0. \emph{Introduction}. On a remarqué, depuis \struck{longtemps}
l'invention de la cohomologie $\ell$-adique [ ], que pour une variété
algébrique $X$ définie sur un corps $k$, de clôture algébrique $\overline{k}$,
les opérations naturelles du groupe de Galois $\pi = \operatorname{Gal}(\overline{k}/k)$
sur la cohomologie $\ell$-adique $H(X_{\overline{k}}, \mathbb{Z}_{\ell})$ de la
variété « géométrique » $X_{\overline{k}}$ définissait sur cette
\struck{dernière} cohomologie une structure supplémentaire remarquable,
donnant lieu notamment aux remarquables conjectures de Tate [ ]. Il semble
probable que ces dernières soient nettement plus éloignées de
\struck{solution} leur solution que les conjectures de Weil [ , ], qui, comme
on sait, peuvent s'interpréter également en termes d'actions galoisiennes sur
des groupes de cohomologie $\ell$-adiques. Notre propos dans le présent
travail est de montrer comment on peut, des conjectures de Weil et (le cas
échéant) de la résolution des singularités, tirer des conséquences de nature
« géométrique » en apparence, i.e.\ relatives à un corps de base
algébriquement clos. Le principe de cette méthode, qui consiste donc
essentiellement à ramener un énoncé géométrique au cas d'un corps de base
fini, est expliqué au N° 1, et des exemples d'application de cette méthode
sont donnés dans les numéros 2 à 7. L'idée de ces raisonnements, liés à
la considération systématique des « poids » des modules galoisiens
provenant de la cohomologie $\ell$-adique, remonte à 1964 ; des exemples
analogues \struck{même principe} se trouvent développés dans SGA 7 (I
et II) ; dans le séminaire oral était développé également l'exemple du
n° 2 du présent travail.
\note{les ajouts « l'invention de la cohomologie $\ell$-adique [ ] »,
« cohomologie », « la », « systématique », « analogues », « développé » et
« sujette à caution » sont de sa main ; une parenthèse dactylographiée
« (qui remonte à 1964) », placée après « relatives à un corps », est biffée à
la machine}

Bien entendu, la portée des « résultats » de notre travail est
nécessairement \struck{limitée} sujette à caution, puisqu'ils
dépendent de conjectures non prouvées à l'heure actuelle, savoir les
conjectures de Weil et la résolution des singularités (sous la forme forte de
Hironaka) en caractéristique

\page{97}
quelconque. Par expérience personnelle, nous pensons cependant que la
compréhension \struck{heuristique quoique pour le moment largement
heuristique}, obtenue grâce aux arguments de poids et de niveaux, quoiqu'elle
soit pour l'instant entièrement heuristique, peut être la source de la
découverte de propriétés importantes et inattendues pour la cohomologie des
variétés algébriques, et qu'elle mérite pour cette raison de ne pas être
dédaignée. De plus, dans le cas des corps de caractéristique nulle,
\struck{des travaux récents fort remarquables les magnifiques} le récent et
magnifique travail de P.\ DELIGNE sur la théorie de Hodge généralisée
des variétés algébriques quelconques sur le corps $\mathbb{C}$ des nombres
complexes, justifie dès maintenant par voie transcendante le « yoga »
des poids et niveaux dans le cadre de la cohomologie de De Rham (tout en
ouvrant à la théorie cohomologique en caractéristique nulle
\struck{des variétés algébriques} des voies entièrement nouvelles). Il
\struck{est} semble probable d'autre part que pour le cas des
\struck{Pour terminer cette introduction} variétés algébriques sur les corps
de caractéristique non nulle, une justification de ce même yoga doive
attendre la \struck{résolution} preuve des conjectures de Weil (ou
mieux, des « conjectures standart » de [ ]).
\note{les ajouts de ce feuillet sont dactylographiés en interligne, sauf les
soulignements et quelques lettres, de sa main}

Pour \struck{signaler} terminer cette introduction, signalons que les
filtrations par poids et niveaux des espaces de cohomologie $\ell$-adiques,
pour $\ell$ variable, sont certainement reliées entre elles de la façon
habituelle, et doivent donner \struck{donnant} en particulier des
espaces vectoriels sur $\mathbb{Q}_{\ell}$ dont le rang est indépendant de
$\ell$. En termes de la théorie conjecturale des motifs (suspendue entre
autres aux « conjectures standart » déjà mentionnées), ceci devrait
s'exprimer par le fait que les structures de poids et de niveau des
cohomologies $\ell$-adiques se placent déjà \struck{se placent} au niveau de
\struck{motifs déjà} la \emph{cohomologie} \struck{des} \emph{motivique},
\struck{\ill{} Dans la} Ce point de vue est en tous cas en accord avec
la théorie de Deligne, \struck{car} puisque sa « filtration par le poids » de
la cohomologie de De Rham provient d'une filtration de la cohomologie à
coefficients \emph{rationnels}. \struck{(et non seulement}

\page{98}
\struck{Introduction 1.}
\note{en haut à gauche, biffé de hachures}

\textbf{Filtration et poids \struck{pour} des espaces de cohomologie des
variétés algébriques}
\note{titre souligné, de sa main}

\begin{enumerate}
  \item[0.] Introduction
  \item[1.] Notion de \emph{poids} et notion de \emph{niveau} pour un
  module \uncertain{pro}-galoisien sur un corps. \struck{\ill{}}
  \item[2.] Étude d'un schéma privé d'un diviseur à croisements normaux.
  \item[3.] Poids et niveaux de $H^{i}_{\uncertain{*}}(X)$ ($X$ lisse)
  \item[4.] Poids et niveaux de $H^{i}_{!}(X)$ ($X$ séparé).
  \struck{\ill{}}
  \item[5.] Poids et niveaux des $H^{i}(X)$ ($X$ quelconque)
  \item[6.] Cas $X$ propre : les plus grands \uncertain{quotients} de
  $H^{i}(X)$ de poids $i$, et résolution des singularités.
  \struck{Applications.}
  \item[7.] Cas $X$ propre et lisse : filtration de $X$ par le niveau.
  \item[8.] Nombres de Betti virtuels de Serre.
  \item[8.] Contre-exemples divers.
  \item[9.] Cas des coefficients tordus ---
\end{enumerate}
\note{les numéros 5, 7 et 8 sont surchargés : un « 8 » est récrit en « 7 »
devant « Nombres de Betti virtuels », et le « 8 » en gras de « Contre-exemples
divers » recouvre un autre chiffre ; on transcrit la numérotation telle
qu'elle se lit, avec son doublon}

[*] SGA 4

[ ] P.\ Deligne / Hodge

[ ] Tate / conjectures

[ ] Kleiman conj.\ standart

[ ] Weil conj.

[ ] Deligne--Grothendieck \emph{SGA 7}

\page{99}
\note{feuillet dactylographié en anglais, dont l'auteur n'est pas nommé sur la
page ; les formules et quelques corrections interlinéaires sont ajoutées à la
main, d'une main qui n'est pas sûrement la sienne. Il se rattache au renvoi
« Kleiman / conj.\ standart » de la p.~98}

\textbf{2. Formalism of Algebraic Cohomology Classes.}

\emph{Cohomology Class Associated to a Cycle}. Let $C^{*}(X)$ be the
intersection ring of cycles on $X$ modulo rational equivalence,
tensored by $\mathbb{Q}$ (the Chow ring of $X$). Then, for each $i
\leqslant n$, we define a homomorphism
\[
  \rho_{X} : C^{i}(X) \longrightarrow H^{2i}_{\ell}(X)
\]
as follows. Let $Z$ be an irreducible subscheme of codimension $i$. The
$\ell$-adic Poincaré duality theorem asserts that the cup-product pairing
\[
  H^{j}_{\ell}(X) \times H^{2n-j}_{\ell}(X) \longrightarrow H^{2n}_{\ell}(X) \xrightarrow[\ \sim\ ]{\ \kappa_{X}\ } \mathbb{Q}_{\ell} \qquad (j = 0, \ldots, 2n)
\]
is a perfect duality of finite dimensional vector spaces, and so the element
$\rho_{X}(Z) \in H^{2i}_{\ell}(X)$ is uniquely determined by the condition
\[
  \kappa_{X}(\alpha \cdot \rho_{X}(Z)) = \kappa_{Z}(\operatorname{res}_{Z}(\alpha)) \quad \text{for all } \alpha \in H^{2n-2i}_{\ell}(X) .
\]
We then extend $\rho_{X}$ to all of $C^{i}(X)$ by linearity. Further, this
homomorphism takes intersection product into cup-product, and thus it extends
to a homomorphism of « algebras » \struck{on} over $\mathbb{Q}$
\[
  \rho_{X} : C^{*}(X) \longrightarrow H^{*}_{\ell}(X) .
\]
The image of a cycle under this homomorphism is called the \emph{cohomology
class} associated with the cycle, and we denote $\rho_{X}(C^{*}(X))$ by
$C^{*}_{\ell}(X)$. A cycle is defined to be $\mathbb{Q}_{\ell}$-cohomologically
equivalent to $0$ if it is in the kernel of the homomorphism $\rho_{X}$.

\page{100}
\[
  \left\{
  \begin{array}{l}
    X = \mathbb{P}^{2} - T' \\
    Y = T - T' \\
    U = \mathbb{P}^{2} - T
  \end{array}
  \right.
\]
\note{à droite, deux croquis : $T$, deux droites qui se coupent, portant les
points $a$ et $b$, $T' = \{a, b\}$ ; et $\widetilde{T}$, les deux droites
séparées, portant l'une $\{a\}$, l'autre $\{b\}$}

\begin{tikzcd}[column sep=small]
  H^{1}(\widetilde{Y}) \arrow[dr] & & \\
  H^{3}_{Y}(X) \arrow[r] & H^{3}(X) \arrow[r, no head] & H^{3}(U)
\end{tikzcd}

\note{sous le $H^{3}$ du milieu, un signe qui pourrait être l'indice $!$ ;
l'exposant de $H^{1}(\widetilde{Y})$ est surchargé}

\begin{tikzcd}[column sep=small]
  H^{1}_{!}(Y) & H^{1}_{!}(X) \arrow[l] \arrow[dl] & H^{1}_{!}(U) \arrow[l] \\
  H^{1}_{!}(\widetilde{Y}) & &
\end{tikzcd}

avec $H^{1}_{!}(U) = 0$ car $T$ connexe.

\begin{tikzcd}[column sep=small]
  H^{0}(\widetilde{T}') \arrow[r] & H^{0}(\widetilde{T}) \arrow[r] & H^{1}_{!}(\widetilde{Y}) \arrow[r] & H^{1}(\widetilde{T}') \\
  H^{0}(\mathbb{P}^{2}) \arrow[r] \arrow[u] & H^{0}(T') \arrow[r] \arrow[u] & H^{1}_{!}(X) \arrow[r] \arrow[u] & H^{1}(\mathbb{P}^{2}) \arrow[u] \\
  H^{0}(\mathbb{P}^{2}) \arrow[r] & H^{0}(T) \arrow[r] \arrow[u] & H^{1}_{!}(U) \arrow[r] \arrow[u] & H^{1}(\mathbb{P}^{2}) \arrow[u, no head, "\Vert" description]
\end{tikzcd}

\note{sous le dernier $H^{1}(\mathbb{P}^{2})$, un second signe d'égalité
vertical, sans rien dessous}

\marginal{$T$ \emph{connexe}, $T'$ disconnexe, \struck{et} contenu dans
\uncertain{$\operatorname{Sing} T$}, \ill{}, et toute comp.\ de $T'$
\uncertain{contient} au plus un point de $T'$.}

\begin{tikzcd}[column sep=small]
  & H^{0}(T')/\operatorname{Im} H^{0}(\mathbb{P}^{2}) \arrow[dl, "\text{inj ?}"'] & H^{0}(T)/\operatorname{Im} H^{0}(\mathbb{P}^{2}) \arrow[l] \\
  H^{0}(\widetilde{T})/\operatorname{Im} H^{0}(\widetilde{T}') & &
\end{tikzcd}

avec $H^{0}(T)/\operatorname{Im} H^{0}(\mathbb{P}^{2}) = 0$ ; et, à gauche de
$H^{0}(\widetilde{T})/\operatorname{Im} H^{0}(\widetilde{T}')$, « $H^{1}_{!}(\widetilde{Y})
\longrightarrow$ », sous lequel \uncertain{$\neq 0$}.
\note{la flèche issue de $H^{1}_{!}(\widetilde{Y})$ et son sens sont
incertains ; les égalités à $0$ sont écrites sous les termes, avec un signe
d'égalité vertical}

\end{document}
